\textbf{Model.} With the gyro measurement $q_m$ (pitch rate, column 2 of the gyro) used as a known input,
\begin{equation}
\dot\theta = q_m + w,\qquad y = \theta_m + v,\qquad w\sim(0,Q),\; v\sim(0,R),
\end{equation}
where $\theta_m=\operatorname{atan2}\!\big(a_x,\sqrt{a_y^2+a_z^2}\big)$ is the raw inclinometer. Discretised at $T_s$ (mean spacing of the
fresh navdata packets, $T_s=13.1$\,ms $\approx 1/76.5$\,s):
$\theta[k{+}1]=\theta[k]+T_s q_m[k]+T_s w[k]$, $y[k]=\theta[k]+v[k]$, i.e.\ $A=1$, $B=T_s$, $G=T_s$, $C=1$.
The steady-state gain comes from \texttt{kalman(ss(1,[Ts Ts],1,[0 0],Ts),Q,R)}
(Control Toolbox, available in R2025b), which returns the innovation gain $M$ of the current estimator:
$\hat\theta[k]=\bar\theta[k]+M\,(y[k]-\bar\theta[k])$, $\bar\theta[k]=\hat\theta[k{-}1]+T_s q_m[k{-}1]$.

\textbf{Choice of $Q$ and $R$} (from the covariances of Sec.~3, \texttt{kf\_pitch\_a.m}).
\emph{Gyro:} $\sigma_q = 0.0007$\,rad/s with motors off (Exp.~A) but $0.0275$\,rad/s in hover (Exp.~C): vibration dominates,
so the in-flight value is used, $Q=\sigma_q^2=7.6\times10^{-4}\,\text{rad}^2/\text{s}^2$.
\emph{Inclinometer:} at rest $\sigma_{\theta_m}=0.34^\circ$ (Exp.~A), but in flight the accelerometer measures specific force
(rotor drag, accelerations), so the error is far larger than the sensor noise and correlates poorly with the attitude. The effective
measurement error was therefore taken as the standard deviation of $\theta_m-\theta_{\rm onboard}$ over the flight segments of Exp.~C and D
($1.64^\circ$ with the mean offset removed), $R=(0.0285\,\text{rad})^2=8.1\times10^{-4}\,\text{rad}^2$. This gives $Q/R=0.93\,\mathrm{s^{-2}}$ and
\begin{equation}
M=0.0125,\qquad \text{closed-loop pole } 1-M=0.9875,\qquad k\simeq\sqrt{Q/R}=0.965\ \text{rad/s}\ (f_c=0.154\ \text{Hz}),
\end{equation}
with a posterior standard deviation of $0.18^\circ$ predicted by the Riccati solution. Using the rest value of $R$ would give $k\approx 4.6$\,rad/s and would
pass most of the accelerometer's flight error into $\hat\theta$.

\textbf{Evaluation on real data.} The filter was run offline on the fresh packets, propagating with the real $\Delta t$ between packets
(Fig.~\ref{fig:q41exp}). Since the onboard estimate $\theta_{\rm onboard}$ is itself an estimate (and differs from the raw inclinometer by a
$0.5$--$0.9^\circ$ offset), the table gives the RMS deviation from it after removing the mean.
\begin{center}\small
\begin{tabular}{lccc}\toprule
 & raw $\theta_m$ & KF & std of KF output / onboard\\\midrule
Exp.~C (hover, 8--28 s) & $0.51^\circ$ & $0.18^\circ$ & $0.37^\circ$ / $0.34^\circ$\\
Exp.~D (pulse, 5--18.4 s) & $2.52^\circ$ & $0.71^\circ$ & $1.62^\circ$ / $2.14^\circ$\\\bottomrule
\end{tabular}
\end{center}
The KF removes most of the high-frequency content of $\theta_m$ (in hover it follows the onboard estimate to $0.18^\circ$) and, in Exp.~D, it recovers the
$\approx\!4^\circ$ pitch pulse between 10 and 13 s even though the raw inclinometer is distorted there (it rises late and decays slowly, because
it is driven by the acceleration produced by the manoeuvre). The KF peak ($\approx 3.8^\circ$) is a little below the onboard one ($\approx 4.5^\circ$) and the filter
undershoots after the pulse: with $k\approx1$\,rad/s the correction from the (wrong) inclinometer is slow but not negligible, so the result
lies between the gyro integral and $\theta_m$.

\begin{figure}[H]\centering
\includegraphics[width=.72\linewidth]{figs/q4_1_exp.pdf}
\caption{Pitch from the raw inclinometer, the onboard estimator and the steady-state KF (Exp.~C and D).}\label{fig:q41exp}
\end{figure}

\textbf{Simulation.} With a synthetic truth (two smooth $8^\circ$ pulses plus a $0.3$\,Hz, $2^\circ$ oscillation), gyro noise $\sqrt Q$ and inclinometer noise $\sqrt R$
(Fig.~\ref{fig:q41sim}), the RMS errors are: raw inclinometer $1.64^\circ$, open-loop gyro integration $0.61^\circ$ (random walk, grows with time), KF $0.23^\circ$ (theory: $0.18^\circ$ in steady state; the extra error is mainly the lag
of the filter on the manoeuvres, which the stationary noise-only Riccati solution does not include).

\begin{figure}[H]\centering
\includegraphics[width=.72\linewidth]{figs/q4_1_sim.pdf}
\caption{Synthetic test: truth, raw measurement, KF and pure gyro integration (top); estimation errors (bottom).}\label{fig:q41sim}
\end{figure}
